\chapter{Tecnologias}
\label{chap2}

\section{Tecnologias de desenvolvimento backend}

\subsection{Java 17}

O Java 17 é uma linguagem de programação orientada a objetos amplamente utilizada no desenvolvimento de sistemas corporativos e aplicações de grande porte. Por se tratar de uma versão de suporte de longo prazo (\textit{Long-Term Support} -- LTS), oferece maior estabilidade, segurança e compatibilidade para projetos que demandam manutenção contínua. Sua adoção contribui para a construção de aplicações robustas, portáveis e com amplo suporte do ecossistema de desenvolvimento.

\subsection{Spring Boot 3}

O Spring Boot 3 é um framework voltado à simplificação do desenvolvimento de aplicações baseadas na plataforma Spring. Sua principal proposta é reduzir a complexidade de configuração inicial, permitindo a criação de sistemas com maior agilidade e padronização. Além disso, fornece mecanismos que facilitam a organização modular da aplicação, a integração entre componentes e a configuração automatizada de diversos recursos.

\subsection{Spring Web}

O Spring Web é o módulo responsável pelo desenvolvimento da camada de comunicação da aplicação, especialmente no que se refere à construção de APIs e ao tratamento de requisições e respostas. Ele possibilita a definição de controladores, rotas e mecanismos de entrada e saída de dados, servindo como base para a interação entre clientes e serviços do sistema.

\subsection{Spring Data JPA / Hibernate}

O Spring Data JPA é uma tecnologia que simplifica o acesso e a manipulação de dados em aplicações Java, oferecendo uma abstração para operações de persistência. Em conjunto com o Hibernate, que atua como implementação do mapeamento objeto-relacional, permite relacionar entidades do sistema com tabelas do banco de dados. Essa combinação reduz a necessidade de escrita manual de consultas básicas, favorecendo produtividade, organização e manutenção do código.

\subsection{Spring Security}

O Spring Security é um framework voltado à segurança de aplicações, oferecendo mecanismos para autenticação, autorização e proteção de recursos. Seu uso permite controlar o acesso às funcionalidades do sistema com base em perfis, permissões e regras previamente definidas. Dessa forma, contribui para reforçar a proteção da aplicação contra acessos indevidos e para estruturar de maneira consistente as políticas de segurança adotadas.

\subsection{JWT}

O JWT (\textit{JSON Web Token}) é um padrão utilizado para representação segura de informações entre partes, sendo amplamente empregado em processos de autenticação e autorização. No contexto de sistemas distribuídos, ele permite que informações sobre a identidade do usuário e suas permissões sejam transmitidas de forma assinada, favorecendo a verificação de acesso sem a necessidade de armazenamento contínuo de sessão no servidor.

\subsection{PostgreSQL}

O PostgreSQL é um sistema gerenciador de banco de dados relacional de código aberto, reconhecido por sua confiabilidade, desempenho e conformidade com padrões SQL. Ele oferece recursos avançados para armazenamento, consulta, integridade e segurança dos dados, sendo amplamente utilizado em aplicações que exigem consistência e escalabilidade na persistência das informações.

\subsection{OpenAPI / Swagger}

O OpenAPI é uma especificação voltada à descrição padronizada de interfaces de programação de aplicações (\textit{Application Programming Interfaces} -- APIs). O Swagger, por sua vez, corresponde a um conjunto de ferramentas que possibilita visualizar, documentar e testar APIs com base nessa especificação. O uso dessas tecnologias favorece a clareza na definição dos contratos da aplicação, facilita a comunicação entre equipes e contribui para a validação e o consumo dos serviços disponibilizados.

\subsection{Maven}

O Maven é uma ferramenta de automação e gerenciamento de projetos utilizada principalmente em aplicações Java. Sua função envolve o controle de dependências, a padronização da estrutura do projeto e a automação de tarefas como compilação, testes e empacotamento. Dessa maneira, contribui para a organização do processo de desenvolvimento e para a reprodutibilidade da construção do software.

\subsection{JUnit / Mockito}

O JUnit é um framework amplamente utilizado para a elaboração e execução de testes automatizados em aplicações Java, especialmente testes unitários. O Mockito, por sua vez, complementa esse processo ao permitir a simulação de comportamentos de objetos e dependências externas. Em conjunto, essas ferramentas contribuem para a verificação do funcionamento do código, para o isolamento de componentes durante os testes e para o aumento da confiabilidade da aplicação.

\subsection{pgAdmin}

O pgAdmin é uma ferramenta gráfica de administração e gerenciamento do PostgreSQL. Sua utilização facilita atividades como visualização de tabelas, execução de consultas, monitoramento do banco e administração de estruturas de dados. Dessa forma, serve como apoio ao desenvolvimento e à manutenção do banco de dados, tornando mais acessíveis operações que poderiam ser realizadas apenas por linha de comando.

\section{Tecnologias de desenvolvimento frontend}

\subsection{React 19}

O React 19 é uma biblioteca JavaScript voltada à construção de interfaces de usuário baseadas em componentes reutilizáveis. Sua abordagem permite dividir a interface em partes independentes, facilitando a organização do código, a manutenção e a evolução da aplicação. Além disso, o React favorece a criação de interfaces dinâmicas e interativas, adequadas ao desenvolvimento de sistemas modernos.

\subsection{TypeScript}

O TypeScript é uma linguagem que estende o JavaScript por meio da adição de tipagem estática e outros recursos voltados à maior segurança no desenvolvimento. Sua utilização contribui para a detecção antecipada de erros, para a melhor documentação do código e para a maior previsibilidade no comportamento da aplicação. Dessa forma, favorece a manutenção e a escalabilidade de projetos de maior porte.

\subsection{Vite}

O Vite é uma ferramenta de construção e execução de aplicações frontend que se destaca pela rapidez no ambiente de desenvolvimento e pela simplicidade de configuração. Ele oferece suporte eficiente ao empacotamento do projeto, à atualização instantânea durante a codificação e à preparação da aplicação para produção. Seu uso contribui para tornar o processo de desenvolvimento mais ágil e produtivo.

\subsection{Material UI (MUI)}

O Material UI (MUI) é uma biblioteca de componentes visuais para aplicações React, baseada nos princípios do \textit{Material Design}. Ela oferece um conjunto de elementos de interface prontos, como botões, tabelas, campos de formulário, menus e diálogos, permitindo a construção de interfaces consistentes, responsivas e visualmente padronizadas. Seu uso contribui para acelerar o desenvolvimento e melhorar a usabilidade da aplicação.

\subsection{Emotion}

O Emotion é uma biblioteca voltada à estilização de interfaces em aplicações JavaScript e React. Ela permite definir estilos diretamente nos componentes, com flexibilidade para composição, reutilização e manutenção do design da aplicação. Essa abordagem favorece a organização dos estilos e a integração entre lógica e apresentação no desenvolvimento da interface.

\subsection{Redux Toolkit}

O Redux Toolkit é uma ferramenta destinada ao gerenciamento de estado global em aplicações frontend. Sua principal finalidade é simplificar o uso da arquitetura Redux, reduzindo a complexidade de configuração e de manipulação do estado compartilhado entre diferentes partes da aplicação. Dessa forma, contribui para a centralização das informações, para a previsibilidade do fluxo de dados e para a melhor organização do código.

\subsection{React Redux}

O React Redux é a biblioteca responsável por integrar o Redux às aplicações desenvolvidas com React. Ela fornece mecanismos para conectar componentes ao estado global da aplicação e para despachar ações que atualizam os dados compartilhados. Seu uso permite que a interface reaja de forma consistente às mudanças de estado, favorecendo a sincronização entre diferentes componentes.

\subsection{React Router DOM}

O React Router DOM é uma biblioteca utilizada para gerenciar a navegação entre diferentes páginas ou visualizações em aplicações React. Ela possibilita definir rotas, associar componentes a caminhos específicos e controlar a transição entre telas sem a necessidade de recarregamento completo da página. Dessa forma, contribui para a organização da navegação e para uma experiência de uso mais fluida.

\subsection{Axios}

O Axios é uma biblioteca utilizada para realizar requisições HTTP em aplicações JavaScript. No contexto do frontend, sua principal função é possibilitar a comunicação com serviços externos, como APIs responsáveis pelo fornecimento e recebimento de dados. Seu uso facilita o envio de requisições, o tratamento de respostas e o gerenciamento de erros na comunicação entre a interface e o backend.

\subsection{ESLint}

O ESLint é uma ferramenta de análise estática de código voltada à identificação de problemas de sintaxe, inconsistências e más práticas em projetos JavaScript e TypeScript. Sua utilização contribui para a padronização do código-fonte, para a melhoria da legibilidade e para a prevenção de erros durante o desenvolvimento. Assim, auxilia na manutenção da qualidade e da consistência do projeto.

\section{Tecnologias de Controle de Versão}


\subsection{Git}

O Git é um sistema de controle de versão distribuído amplamente utilizado no desenvolvimento de software para registrar, organizar e acompanhar as alterações realizadas no código-fonte ao longo do tempo. Por meio dele, torna-se possível manter um histórico detalhado das modificações, recuperar versões anteriores, criar ramificações independentes para desenvolvimento de funcionalidades e integrar mudanças de diferentes colaboradores de forma controlada. Sua utilização contribui para a rastreabilidade das alterações, para a segurança no processo de desenvolvimento e para a manutenção evolutiva do software.

\subsection{GitLab}

O GitLab é uma plataforma de hospedagem e gerenciamento de repositórios baseada no Git, que oferece recursos voltados ao desenvolvimento colaborativo e à automação de processos. Além de permitir o armazenamento centralizado do código-fonte, a ferramenta disponibiliza mecanismos para controle de acesso, revisão de código, gerenciamento de \textit{issues}, acompanhamento de atividades e integração com fluxos de integração contínua e entrega contínua. Dessa forma, o GitLab amplia as possibilidades de organização do projeto, de colaboração entre os membros da equipe e de controle do ciclo de desenvolvimento de software.

\section{Tecnologias de Conteinerização e Distribuição}

\subsection{Docker}

O Docker é uma plataforma voltada à conteinerização de aplicações, permitindo empacotar o software e suas dependências em unidades isoladas de execução, denominadas contêineres. No contexto de integração contínua e entrega contínua, seu uso contribui para a padronização dos ambientes de desenvolvimento, teste e implantação, reduzindo inconsistências entre diferentes etapas do ciclo de vida do software. Além disso, o Docker favorece a reprodutibilidade, a portabilidade e a escalabilidade da aplicação, tornando o processo de construção, validação e disponibilização de novas versões mais previsível e controlado.

\subsection{GitLab CI/CD}
O GitLab CI/CD corresponde ao conjunto de mecanismos de automação disponibilizados pela plataforma GitLab para apoiar práticas de integração contínua e entrega contínua. Por meio dessa abordagem, é possível definir fluxos automatizados para executar etapas como compilação, testes, validações e processos de implantação sempre que alterações são realizadas no repositório. Sua utilização contribui para reduzir atividades manuais, aumentar a confiabilidade do processo de desenvolvimento e permitir que a equipe acompanhe de forma mais sistemática a qualidade e a evolução do software. Dessa forma, o GitLab Actions fortalece a automação do ciclo de desenvolvimento e favorece a disponibilização mais rápida e consistente de novas versões da aplicação.